\section*{Capítulo \ref{ch:b3:conservation-biology} --- Biología de la conservación}

\begin{solution}{exo:b3:conservation-biology:1}
\omterm{def:b3:conservation-biology:small}{Estocasticidad demográfica}: el azar en los nacimientos, las muertes y
los sexos de unos pocos individuos ---el kakapo con cincuenta y uno,
con más machos que hembras. \omterm{def:b3:conservation-biology:small}{Estocasticidad ambiental}: un mal año que
golpea a la vez a todos los individuos ---una sequía, una enfermedad, un
invierno duro que cae sobre una población de unos pocos cientos. \omterm{def:b3:conservation-biology:small}{Efecto
Allee}: el crecimiento per cápita que cae a densidad baja ---la paloma
migratoria, cuyas colonias no podían criar una vez raleadas y cuyos
depredadores restantes ya no quedaban saturados.
\end{solution}

\begin{solution}{exo:b3:conservation-biology:2}
El tamaño de una población ideal ---constante, de apareamiento
aleatorio, con sexos iguales y tamaños de familia de Poisson--- que
perdería heterocigosidad al mismo ritmo que la real. Números desiguales
de machos y hembras reproductores; fluctuaciones de tamaño, que pesan
más en las generaciones menores; varianza del tamaño de familia por
encima de la de Poisson, con unos pocos individuos que producen la
mayoría de las crías (y también el apareamiento no aleatorio y las
generaciones solapadas).
\end{solution}

\begin{solution}{exo:b3:conservation-biology:3}
$S = cA^{z}$. Fracción conservada $= (A'/A)^{z}$: $0.5^{0.25} = 0.84$;
$0.1^{0.25} = 0.56$; $0.01^{0.25} = 0.32$ ---una centésima parte del
área conserva aún, con el tiempo, un tercio de las especies.
\end{solution}

\begin{solution}{exo:b3:conservation-biology:4}
$N_{e} \ge 50$ mantiene la subida de la endogamia por debajo de $1/100$
por generación, el ritmo al que los criadores de animales encuentran
tolerable la depresión: protege de la pérdida de eficacia a corto
plazo. $N_{e} \ge
500$ equilibra la pérdida de variación por deriva con su aporte por
mutación, de modo que la población conserva la variación cuantitativa
que necesita para adaptarse: protege el largo plazo. Las poblaciones
reales necesitan de cinco a diez veces esas cifras en el censo.
\end{solution}

\begin{solution}{exo:b3:conservation-biology:5}
$N_{e} = 4\times 8\times 120/128 = 30$; la heterocigosidad cae $1/60 =
\qty{1.7}{\%}$ por generación; se pierde una cuarta parte cuando $(1 - 1/60)^{t} =
0.75$, con $t = \ln 0.75/\ln(59/60) = 17$ generaciones.
\end{solution}

\begin{solution}{exo:b3:conservation-biology:6}
$1/N_{e} = \tfrac{1}{5}(4/1000 + 1/20) = 0.0108$, $N_{e} = 93$: una sola
generación de veinte ha hecho que cinco generaciones se comporten como
noventa y tres. Heterocigosidad retenida, $(1 - 1/186)^{5} = 0.973$,
frente a $(1 -
1/2000)^{5} = 0.9975$ con un millar constante.
\end{solution}

\begin{solution}{exo:b3:conservation-biology:7}
$F = 1 - (1 - 1/20)^{5} = 0.23$; la eficacia baja
$5\times 2.3 = \qty{11}{\%}$. Toda descendencia de una hembra de Texas y
un macho de Florida tiene $F = 0$; si las ocho recién llegadas son dos
quintas partes de las hembras reproductoras, dos quintas partes de la
generación siguiente son exogámicas y el $F$ medio cae hasta alrededor
de $0.6\times 0.23 = 0.14$ en una generación ---que es lo que mostró la
recuperación.
\end{solution}

\begin{solution}{exo:b3:conservation-biology:8}
$r = \ln(100/31)/8 = \qty{0.15}{yr^{-1}}$. Continuado hasta 2020:
$100\,\mathrm{e}^{0.15\times 17} \approx 1200$. El parque albergó unos
cien: los territorios se llenaron, los uapitíes declinaron, las manadas
mataron a miembros de las otras y se extendieron la sarna y el moquillo
---la dependencia de la densidad se instaló una vez ocupado el terreno
vacío.
\end{solution}

\begin{solution}{exo:b3:conservation-biology:9}
(a) Un declive del \qty{60}{\%} en diez años supera el \qty{50}{\%}: En
Peligro. (b) $180 < 250$ individuos maduros: En Peligro Crítico. (c)
Área de distribución de \qty{3000}{km^2}, por debajo de \qty{5000}{}: En
Peligro. (d) \num{30000} individuos y un declive del \qty{10}{\%}: no se
alcanza ninguno de los umbrales ---no amenazada (a lo sumo Casi
Amenazada, a vigilar).
\end{solution}

\begin{solution}{exo:b3:conservation-biology:10}
Las $N$ parejas fallan todas con probabilidad $(1/2)^{N}$: $0.25$,
$0.031$, $0.001$ y $10^{-6}$. El riesgo demográfico cae
exponencialmente con $N$ y es despreciable más allá de unas pocas
docenas de individuos; por encima de eso, el riesgo que queda es
ambiental, que no cae tan deprisa, porque golpea a todos los individuos
a la vez.
\end{solution}

\begin{solution}{exo:b3:conservation-biology:11}
Un fragmento de una décima parte: $0.1^{0.25} = 0.56$ de las especies
originales. Diez fragmentos albergan entre $0.56$ (si todos albergan
las mismas especies) y, en principio, más que el conjunto (si cada uno
albergara especies distintas ---imposible más allá del acervo
regional). Dónde caigan depende del recambio entre fragmentos: hábitats
distintos en fragmentos distintos empujan el total hacia arriba; las
especies que necesitan áreas grandes o hábitat interior, o que se
pierden de todos los fragmentos por los mismos \omterm{prop:b3:conservation-biology:area}{efectos de borde}, lo
empujan hacia abajo; y la deuda de extinción de cada fragmento se paga
por separado, de modo que el total a largo plazo suele quedar por
debajo del del conjunto.
\end{solution}

\begin{solution}{exo:b3:conservation-biology:12}
La deriva retira heterocigosidad a un ritmo $1/2N_{e}$ por generación;
los inmigrantes reemplazan una fracción $m$ del acervo génico por
alelos tomados de otra parte. Con $mN_{e} \approx 1$, $m \approx 1/N_{e}$,
el doble del ritmo de pérdida, de modo que la variación se repone más
deprisa de lo que decae y la población se mantiene cerca de la fuente
en frecuencias alélicas ($F_{ST} = 1/(1
+ 4N_{e}m) \approx 0.2$). El coste: los alelos favorecidos localmente
se diluyen en $m$ por generación, de modo que la adaptación local solo
sobrevive para los rasgos cuyo coeficiente de selección supera
aproximadamente $1/N_{e}$; las diferencias locales débilmente
seleccionadas quedan anegadas.
\end{solution}

\begin{solution}{pb:b3:conservation-biology:1}
\textbf{1.} $N_{e} = 4\times 20\times 40/60 = 53$, frente a un censo de 60.
\textbf{2.} $1/N_{e} = \tfrac{1}{5}(1/500 + 1/100 + 1/60 + 1/51 + 1/60) =
0.0130$, $N_{e} = 77$; dominan las generaciones de cuello de botella de
51 y 60, y la de 500 apenas cuenta.
\textbf{3.} $F = 1 - (1 - 1/154)^{5} = 0.032$; lineal, $5/154 = 0.032$.
\textbf{4.} $6\times 0.32 = \qty{1.9}{\%}$.
\textbf{5.} $(1 - 1/154)^{5} = 0.968$: se retiene el \qty{96.8}{\%}.
\textbf{6.} Diez generaciones con $N_{e} = 53$: incremento $1 - (1 -
1/107)^{10} = 0.090$, total $F = 1 - 0.968\times 0.910 = 0.12$; pérdida
de eficacia $6\times 1.2 = \qty{7}{\%}$; un siglo.
\textbf{7.} Las parejas de nativo con nativo son $(60/70)^{2} = 0.73$ de
los apareamientos y su descendencia conserva $F \approx 0.03$ (más un
pequeño incremento); todas las demás tienen $F = 0$: en media,
$F \approx 0.73\times 0.035 = 0.026$, una caída de la cuarta parte en
una generación.
\textbf{8.} $F$ es identidad por descendencia dentro de un individuo, y
cualquier descendiente exogámico lo devuelve a cero; la
heterocigosidad es una propiedad de las frecuencias alélicas, y los
alelos de los migrantes se extienden solo a medida que se transmiten.
En una población todavía pequeña, la deriva volverá a retirarlos en
unas decenas de generaciones, de modo que el rescate ha de repetirse o
la población ha de crecer.
\textbf{9.} $N_{e} = N$ con sexos iguales: 500 adultos. Desde 70 con $r =
0.05$: $\ln(500/70)/0.05 = 39$ años.
\textbf{10.} La población ideal tiene tamaños de familia de Poisson, con
varianza igual a la media de dos; $N_{e} = 4N/(2 + V_{k})$, de modo que
hacer que todas las parejas contribuyan por igual ($V_{k} = 0$) da
$N_{e} = 2N$: ningún linaje se pierde por la suerte del nido.
\textbf{11.} Recójase y consérvese su semen para inseminación
artificial, y úsese en varias hembras sin parentesco a lo largo de los
años ---el coste es que sus alelos podrían acabar dominando un acervo
pequeño, de modo que su parte ha de limitarse a un valor cercano a
$1/N_{e}$---; o criopreservar tejido para una clonación o una
producción de gametos posteriores, al coste del retraso y del riesgo
técnico. No hacer nada pierde para siempre los alelos del linaje.
\textbf{12.} \qty{160}{} millones de dólares por unas 430 aves: unos
\num{370000} por ave. Caro por loro, barato frente a una especie ---y
menos que unos pocos kilómetros de autopista, que es la comparación que
los presupuestos hacen en realidad.
\textbf{13.} $(300/2000)^{0.25} = 0.62$: se conservan unas 112 especies
de 180, con una deuda de extinción de 68.
\textbf{14.} Cada fragmento, $(100/2000)^{0.25}\times 180 = 85$
especies. Enteramente distintas: hasta 255, imposible por encima del
acervo de 180; idénticas: 85. Lo que lo decide es cuán distintos son
los hábitats de los fragmentos y cuántas especies pueden persistir en
\qty{100}{km^2}; de forma realista, entre 85 y 112.
\textbf{15.} $(320/2000)^{0.25}\times 180 = 114$, frente a 85 con tres
fragmentos aislados e idénticos: el corredor deja que la
recolonización y el flujo génico hagan que tres poblaciones pequeñas se
comporten como una mayor.
\textbf{16.} $N_{e} = 50$ con sexos iguales son 25 parejas:
\qty{1250}{km^2}. Los \qty{300}{km^2} enteros albergan seis parejas,
$N_{e} \approx 12$; un fragmento, dos. Ninguno basta: el depredador se
pierde sea cual sea el recuento de especies, y proteger el área que
necesita protege todo lo menor ---una especie paraguas.
\textbf{17.} $\ln(200/10)/0.03 = 100$ años: la deuda se paga a lo largo
de un siglo, mucho después de talado el bosque.
\textbf{18.} Los bordes despojan de hábitat interior a todos los
fragmentos, de modo que el área utilizable es menor que la
cartografiada y la pérdida es mayor que la que predice la ley; una
matriz de bosque secundario o de setos que algunas especies puedan usar
añade área efectiva y hace menor la pérdida.
\textbf{19.} $r = \ln(100/31)/8 = \qty{0.146}{yr^{-1}}$; tiempo de
duplicación $\ln 2/0.146 = 4.7$ años.
\textbf{20.} $N(8) = 120/[1 + (120/31 - 1)\mathrm{e}^{-1.17}] = 120/(1 +
2.87\times 0.31) = 63$. Los 100 observados son el propio valor
exponencial ---$r$ se ajustó con esos mismos dos puntos--- frente a los
63 logísticos: en los primeros años los lobos no encontraron
dependencia de la densidad ---territorios vacíos y uapitíes
abundantes.
\textbf{21.} $\ln(4000/\num{17000})/20 = -0.072$: un \qty{7}{\%} anual.
Cien lobos que se llevan \num{2000} uapitíes al año de una manada de
\num{8000} de media son una cuarta parte al año ---más que suficiente
para impulsar el declive por sí solos, aunque la caza fuera del parque,
la sequía y la depredación de crías por los osos grises lo
compartieran.
\textbf{22.} Sesenta individuos maduros, por debajo de 250: En Peligro
Crítico. El depredador, con $N_{e} \approx 12$ (unas pocas docenas de
animales): también En Peligro Crítico.
\textbf{23.} De $3\times 10^{9}$ a $1$ en 44 años: $r = -\ln(3\times
10^{9})/44 = -0.50$ al año, con una reducción a la mitad cada diecisiete
meses. Ninguna cosecha retira la mitad de miles de millones cada año; la
caída fue más lenta al principio y mucho más rápida al final, porque una
especie de cría colonial por debajo de una densidad umbral deja de
reproducirse mientras sus depredadores ya no quedan saturados ---el
\omterm{def:b3:conservation-biology:small}{efecto Allee} convirtió un declive en un derrumbe.
\textbf{24.} En 1850 la población se contaba por miles de millones, su
tasa de crecimiento medida era positiva y su varianza pequeña, y un
análisis de viabilidad proyecta la dinámica que se le da: habría puesto
el riesgo en cero. Se le habría escapado el forzamiento externo ---los
ferrocarriles y el telégrafo, que permitieron a los cazadores seguir y
vaciar todos los dormideros, y el desmonte de los bosques--- y el umbral
de Allee de una especie de cría colonial, ninguno de los cuales está en
la demografía de una especie abundante; un análisis de viabilidad
extrapola, y las causas de la extinción suelen ser nuevas.
\textbf{25.} $N_{e} \approx 53$ hoy; $F \approx 0.12$ tras diez
generaciones más con 60 individuos (un \qty{7}{\%} de la eficacia); el
bosque encogido conserva alrededor del \qty{62}{\%} de sus especies.
\end{solution}
